\documentclass[11pt,spanish]{article}

% =========================================================
% PLANTILLA BASE PARA INFORMES ACADÉMICOS / TÉCNICOS
% =========================================================

% --- Idioma y codificación ---
\usepackage[spanish,es-tabla]{babel}
\usepackage[T1]{fontenc}
\usepackage{textcomp}
\usepackage{newunicodechar}
\newunicodechar{₡}{\textcolonmonetary}

% --- Márgenes / papel ---
\usepackage[
    left=2.5cm,
    right=2.5cm,
    top=2.5cm,
    bottom=2.5cm,
    headsep=0.5cm,
    footskip=0.8cm
]{geometry}

% --- Matemáticas ---
\usepackage{amsmath}

% --- Tablas, enlaces y elementos visuales ---
\usepackage{array}
\usepackage{tabularx}
\usepackage{booktabs}
\usepackage{longtable}
\usepackage{graphicx}
\usepackage{float}
\usepackage{xcolor}
\usepackage{tikz}
\usetikzlibrary{positioning,shapes.geometric,arrows.meta,calc}
\usepackage{enumitem}
\usepackage{ragged2e}
\usepackage[hidelinks]{hyperref}

% --- Encabezado y pie de página ---
\usepackage{fancyhdr}
\usepackage{lastpage}

% --- Interlineado ---
\linespread{1.2}
\sloppy

% =========================================================
% DATOS DEL DOCUMENTO
% =========================================================
\newcommand{\Institucion}{Instituto Tecnológico de Costa Rica}
\newcommand{\Campus}{Campus Tecnológico Central Cartago}
\newcommand{\DocumentoTitulo}{Registro de Riesgos}
\newcommand{\Curso}{Administración de proyectos}
\newcommand{\Profesor}{Alicia Marcela Salazar Hernández}
\newcommand{\FechaDocumento}{10 de junio de 2026}
\newcommand{\PeriodoAcademico}{I Semestre}
\newcommand{\AnioDocumento}{2026}
\newcommand{\AutoresPortada}{%
    \begin{tabular}{l@{\hspace{1.5cm}}r}
        Mauricio González Prendas & 2024143009 \\
        Isaac Jiménez Blanco & 2024202804 \\
        Tamara Robles Camacho & 2024099342 \\
    \end{tabular}%
}
\newcommand{\DocHeaderLeft}{}
\newcommand{\DocHeaderRight}{Registro de Riesgos}
\newcommand{\DocFooterRight}{Página~\thepage\ de~\pageref{LastPage}}

% Metadatos PDF.
\title{\DocumentoTitulo}
\author{\AutoresPortada}
\date{\FechaDocumento}

% =========================================================
% CONFIGURACIÓN DE TABLAS Y UTILIDADES
% =========================================================
\newcolumntype{Y}{>{\RaggedRight\arraybackslash}X}
\newcolumntype{L}[1]{>{\RaggedRight\arraybackslash}p{#1}}
\renewcommand{\arraystretch}{1.22}
\setlength{\LTpre}{0.5em}
\setlength{\LTpost}{0.5em}

% Imagen segura: si el archivo no existe, muestra un recuadro de marcador.
\newcommand{\SafeImage}[2][0.92\textwidth]{%
    \IfFileExists{#2}{%
        \includegraphics[width=#1]{#2}%
    }{%
        \fbox{%
            \begin{minipage}[c][4cm][c]{#1}
            \centering
            Imagen pendiente:\\
            \texttt{#2}
            \end{minipage}%
        }%
    }%
}

% Figura reutilizable.
\newcommand{\EvidenceFigure}[3]{%
    \begin{figure}[H]
    \centering
    \SafeImage{#1}
    \caption{#2}
    \label{#3}
    \end{figure}%
}

% =========================================================
% ENCABEZADO Y PIE
% =========================================================
\pagestyle{fancy}
\fancyhf{}
\fancyhead[L]{\DocHeaderLeft}
\fancyhead[R]{\DocHeaderRight}
\renewcommand{\headrulewidth}{0.4pt}
\renewcommand{\footrulewidth}{0.4pt}
\fancyfoot[R]{\DocFooterRight}

% Estilo equivalente para páginas horizontales.
\fancypagestyle{widefancy}{%
    \fancyhf{}%
    \fancyhead[L]{\DocHeaderLeft}%
    \fancyhead[R]{\DocHeaderRight}%
    \renewcommand{\headrulewidth}{0.4pt}%
    \renewcommand{\footrulewidth}{0.4pt}%
    \fancyfoot[R]{\DocFooterRight}%
}

% =========================================================
% PÁGINAS HORIZONTALES PARA TABLAS ANCHAS
% =========================================================
\makeatletter
\newcommand{\SyncPageBuilderDimensions}{%
    \global\vsize=\textheight
    \global\@colroom=\textheight
    \global\@colht=\textheight
}
\makeatother

\newenvironment{widepage}{%
    \clearpage
    \hypersetup{pageanchor=false}%
    \paperwidth=297mm
    \paperheight=210mm
    \pdfpagewidth=297mm
    \pdfpageheight=210mm
    \newgeometry{left=2.5cm,right=2.5cm,top=2.5cm,bottom=2.5cm,headsep=0.5cm,footskip=0.8cm}%
    \setlength{\textwidth}{243mm}%
    \setlength{\textheight}{156mm}%
    \setlength{\linewidth}{\textwidth}%
    \setlength{\columnwidth}{\textwidth}%
    \hsize=\textwidth
    \setlength{\headwidth}{\textwidth}%
    \SyncPageBuilderDimensions
    \pagestyle{widefancy}%
}{%
    \clearpage
    \paperwidth=210mm
    \paperheight=297mm
    \pdfpagewidth=210mm
    \pdfpageheight=297mm
    \restoregeometry
    \SyncPageBuilderDimensions
    \hypersetup{pageanchor=true}%
    \pagestyle{fancy}%
}

% =========================================================
% PORTADA
% =========================================================
\newcommand{\MakeCover}{%
\begin{titlepage}
\thispagestyle{empty}
\centering

{\bfseries \huge \Institucion \par}
\vspace{0.25cm}
{\LARGE \Campus \par}

\vfill

{\huge \DocumentoTitulo \par}

\vfill

{\bfseries \LARGE Profesora \par}
{\Large \Profesor \par}

\vfill

{\bfseries \LARGE Estudiantes \par}
{\Large \AutoresPortada \par}

\vfill

{\bfseries\LARGE Fecha \par}
{\Large \FechaDocumento \par}

\vfill

{\LARGE \PeriodoAcademico \par}
{\LARGE \AnioDocumento \par}

\end{titlepage}%
}

% =========================================================
% DOCUMENTO
% =========================================================
\begin{document}
\hypersetup{pageanchor=false}

% =========================
% Portada
% =========================
\MakeCover

% =========================
% Índice
% =========================
\pagenumbering{roman}
\pagestyle{empty}
\addtocontents{toc}{\protect\thispagestyle{empty}}
\tableofcontents
\clearpage
\pagenumbering{arabic}
\hypersetup{pageanchor=true}
\pagestyle{fancy}

% =========================
% Contenido
% =========================
\section{Información general del documento}

\begin{table}[H]
\centering
\begin{tabularx}{\textwidth}{p{5cm} Y}
\toprule
\textbf{Nombre del proyecto} & Plataforma Universitaria Integrada \\
\textbf{Organización} & Universidad Tecnológica La Mejor \\
\textbf{Documento} & \DocumentoTitulo \\
\textbf{Directora del proyecto} & Tamara Robles Camacho \\
\textbf{Fecha} & \FechaDocumento \\
\bottomrule
\end{tabularx}
\end{table}

\section{Introducción}

El presente documento detalla el registro de riesgos elaborado para el proyecto denominado Plataforma Universitaria Integrada, el cual es desarrollado para la Universidad Tecnológica La Mejor. Este registro constituye un componente fundamental dentro de la planificación y control del proyecto, para lo cual se tiene como finalidad principal la identificación de las posibles eventualidades que podrían afectar el alcance, el cronograma, el costo, la calidad, la comunicación, el desarrollo técnico de la interfaz de usuario y el cierre exitoso del proyecto.

La gestión oportuna de estos riesgos es de suma importancia debido a que permite anticipar diversas situaciones adversas que podrían generar atrasos, retrabajo, pérdida de información, defectos en el prototipo final o inconsistencias documentales. En consecuencia, en esta iniciativa en particular, los riesgos deben ser analizados de manera sumamente minuciosa ya que el equipo de trabajo cuenta con un cronograma planificado de tres meses y debe cumplir con múltiples entregables documentales, además de gestionar roles combinados y un prototipo de interfaz navegable cuya finalización está programada según la fecha aprobada en Microsoft Project.

El proyecto consiste propiamente en el desarrollo de una Plataforma Universitaria Integrada implementada a nivel de prototipo web navegable, para lo cual se deben representar de forma interactiva diversos módulos como el inicio de sesión, el panel de control principal, el portal estudiantil, el portal docente, el portal administrativo, el portal financiero y el panel ejecutivo. Asimismo, además de este prototipo técnico, la propuesta incorpora múltiples documentos de gestión orientados a complementar y consolidar la planificación estructurada del proyecto.

El alcance detallado del proyecto establece con claridad que no se desarrollará la lógica funcional del servidor, bases de datos reales, sistemas de autenticación institucional integrados, pasarelas de pago reales, migración de datos históricos ni un despliegue en producción, de manera que varios riesgos identificados se orientan a evitar que tanto el equipo de desarrollo como los interesados esperen funcionalidades que se encuentran fuera del alcance aprobado. Por lo tanto, otros riesgos importantes se asocian directamente con la coherencia de la documentación, la disponibilidad de los miembros del equipo, la calidad técnica del prototipo navegable, la efectividad en la comunicación interna y el uso de las distintas herramientas especializadas de desarrollo.

Este registro de riesgos se ha estructurado a partir del contexto y de los objetivos generales de la iniciativa, de manera que se asegura que las amenazas detectadas no se planteen de forma aislada, sino estrechamente vinculadas con las condiciones reales del entorno y con los elementos organizacionales definidos formalmente durante la etapa de planificación.

\section{Propósito del registro de riesgos}

El propósito central de este registro de riesgos es documentar de forma estructurada los eventos inciertos que posean el potencial de afectar de manera negativa el cumplimiento de las metas del proyecto, para lo cual este documento facilita la identificación precisa de las causas de cada riesgo, su impacto estimado, su nivel de prioridad, el responsable asignado para su monitoreo y la estrategia de respuesta adecuada que se deba aplicar.

Este instrumento se orienta al cumplimiento de objetivos específicos detallados a continuación.

\begin{itemize}[leftmargin=*, label={\textbullet}]
    \item Identificar los riesgos principales del proyecto.
    \item Analizar la causa y el posible impacto de cada riesgo.
    \item Evaluar la probabilidad y el impacto mediante una escala definida.
    \item Clasificar los riesgos según su nivel de criticidad.
    \item Asignar responsables para dar seguimiento a cada riesgo.
    \item Definir acciones de respuesta para prevenir o reducir sus efectos.
    \item Mantener trazabilidad entre riesgos, alcance, cronograma, calidad, comunicaciones y cambios.
    \item Apoyar la toma de decisiones del equipo durante la ejecución y cierre del proyecto.
\end{itemize}

Asimismo, el registro de riesgos de la Plataforma Universitaria Integrada resulta fundamental para priorizar aquellos eventos que demandan una atención inmediata por parte de la gestión, debido a que no todas las amenazas presentan la misma relevancia para el proyecto. Por lo tanto, un error menor de formato en la documentación no representará el mismo impacto que un atraso en la codificación de la interfaz navegable o una modificación extemporánea en el alcance, por lo cual se implementará una evaluación basada en probabilidad e impacto con el fin de determinar las prioridades de atención.

\section{Metodología de evaluación de riesgos}

Para evaluar de manera adecuada las amenazas identificadas se utilizará un enfoque cualitativo centrado en la probabilidad y el impacto, de modo que cada riesgo sea analizado a partir de la posibilidad de su ocurrencia y del efecto directo que causaría sobre el desarrollo si llegara a materializarse.

Esta valoración se ejecutará empleando una escala cuantitativa de uno a cinco tanto para la probabilidad como para el impacto, para lo cual posteriormente se calculará el nivel de riesgo definitivo mediante el producto de ambos factores.

La determinación del nivel de riesgo se realiza mediante la siguiente ecuación.

\begin{center}
Nivel de riesgo = Probabilidad x Impacto
\end{center}

A partir del valor obtenido se procederá a clasificar cada riesgo en las categorías de bajo, medio, alto o crítico, lo que implica que esta segmentación facilitará al equipo la identificación inmediata de las amenazas que requieren un monitoreo prioritario.

El análisis detallado de cada riesgo contemplará los elementos que se enumeran a continuación.

\begin{itemize}[leftmargin=*, label={\textbullet}]
    \item Código del riesgo.
    \item Categoría.
    \item Descripción del riesgo.
    \item Causa.
    \item Impacto en el proyecto.
    \item Probabilidad.
    \item Impacto.
    \item Nivel resultante.
    \item Clasificación.
    \item Responsable.
    \item Estrategia de respuesta.
    \item Acción de respuesta.
    \item Estado.
\end{itemize}

\section{Escala de probabilidad e impacto}

La evaluación de la probabilidad y el impacto en este registro se encuentra alineada con la matriz de probabilidad e impacto definida para la organización, en la cual se detalla la escala de referencia junto con la ubicación de las trece amenazas identificadas. Por lo tanto, con el propósito de asegurar la consistencia entre los documentos del proyecto, este registro adopta de manera estricta los mismos valores de probabilidad, impacto, nivel de riesgo y clasificación acordados en dicho instrumento de análisis.

\section{Estrategias de respuesta ante riesgos}

Para cada riesgo identificado se establecerá una estrategia de respuesta específica, de manera que las acciones correspondan a una de las categorías descritas a continuación.

\subsection{Evitar}

Esta estrategia consiste en modificar los procesos o planes de trabajo con el objetivo de eliminar por completo la probabilidad de ocurrencia de la amenaza, de modo que se aplica prioritariamente en aquellos eventos que podrían comprometer gravemente el alcance, el cronograma o la calidad de los entregables.

\subsection{Mitigar}

Esta alternativa busca reducir la probabilidad de ocurrencia o la severidad del impacto mediante la ejecución de acciones preventivas, lo cual representa la estrategia de mayor uso en el proyecto debido a que la mayoría de los riesgos pueden ser gestionados mediante actividades constantes de revisión, comunicación y seguimiento.

\subsection{Transferir}

Esta modalidad implica trasladar la responsabilidad o el impacto financiero y operativo del riesgo hacia un tercero ajeno al equipo de desarrollo, la cual se empleará de manera limitada en esta iniciativa debido a que solo se aplicará ante dependencias tecnológicas externas o fallos en los servicios de almacenamiento en la nube.

\subsection{Aceptar}

Esta estrategia consiste en asumir formalmente la existencia del riesgo y realizar un monitoreo pasivo sin implementar acciones preventivas inmediatas, de manera que se utiliza en situaciones donde el impacto es mínimo o cuando el costo de mitigación supera los beneficios esperados de la intervención.

\section{Fuentes utilizadas para identificar riesgos}

La identificación de las amenazas se llevó a cabo mediante una revisión exhaustiva de las bases de planificación del proyecto, para lo cual las fuentes documentales utilizadas corresponden a las que se mencionan a continuación.

\begin{itemize}[leftmargin=*, label={\textbullet}]
    \item Caso de negocio.
    \item Acta de constitución del proyecto.
    \item Registro de interesados.
    \item Definición del alcance.
    \item Cronograma del proyecto detallado en Microsoft Project.
    \item Organigrama estructurado del proyecto.
    \item Plan de gestión de calidad.
    \item Plan de gestión de comunicaciones.
    \item Restricciones asociadas a los plazos y al alcance de los entregables.
    \item Distribución interna de roles entre los miembros del equipo.
    \item Requerimientos técnicos para el desarrollo del prototipo de interfaz navegable.
    \item Herramientas tecnológicas de trabajo cooperativo incluyendo GitHub, Microsoft Project, procesadores de texto, plataformas de comunicación y correo electrónico.
\end{itemize}

La consulta detallada de estos recursos propició la detección de amenazas asociadas con los plazos de entrega, el alcance acordado, la calidad de los productos, la coherencia de la documentación, los canales de comunicación, el funcionamiento de las herramientas y la coordinación técnica del desarrollo.

\begin{widepage}
\section{Registro de riesgos del proyecto}
{\tiny
\setlength{\tabcolsep}{2pt}
\begin{longtable}{L{0.9cm} L{1.9cm} L{2.1cm} L{2.4cm} L{2.3cm} L{1.8cm} L{1.2cm} L{0.9cm} L{1.7cm} L{2.2cm} L{1.4cm} L{2.6cm} L{1.0cm}}
\caption{Registro de riesgos del proyecto}\label{tab-registro-riesgos}\\[-0.5em]
\toprule
\textbf{Código} & \textbf{Categoría} & \textbf{Riesgo} & \textbf{Causa} & \textbf{Impacto en el proyecto} & \textbf{Probabilidad} & \textbf{Impacto} & \textbf{Nivel} & \textbf{Clasificación} & \textbf{Responsable} & \textbf{Estrategia} & \textbf{Acción de respuesta} & \textbf{Estado} \\
\midrule
\endfirsthead
\toprule
\textbf{Código} & \textbf{Categoría} & \textbf{Riesgo} & \textbf{Causa} & \textbf{Impacto en el proyecto} & \textbf{Probabilidad} & \textbf{Impacto} & \textbf{Nivel} & \textbf{Clasificación} & \textbf{Responsable} & \textbf{Estrategia} & \textbf{Acción de respuesta} & \textbf{Estado} \\
\midrule
\endhead
R-01 & Cronograma & Restricción fuerte de tiempo & El proyecto debe entregarse según la fecha final establecida en el cronograma aprobado y contiene varios entregables documentales y técnicos & Atraso en documentos, reducción de calidad o entrega incompleta & 4 & 5 & 20 & Crítico & Tamara Robles Camacho & Mitigar & Dar seguimiento frecuente al cronograma, priorizar tareas críticas y redistribuir trabajo si una actividad se atrasa & Abierto \\
R-02 & Alcance & Alcance creciente & Se pueden proponer funcionalidades fuera del alcance, como lógica de servidor, base de datos real o pagos reales & Aumento del trabajo, atraso en la interfaz y contradicción con la definición del alcance & 4 & 5 & 20 & Crítico & Tamara Robles Camacho e Isaac Jiménez Blanco & Evitar & Revisar toda nueva solicitud contra el alcance aprobado y aplicar control formal de cambios & Abierto \\
R-04 & Desarrollo & Subestimación del esfuerzo de la interfaz & El prototipo incluye varios módulos, pantallas y navegación entre portales & El prototipo puede quedar incompleto o con navegación limitada & 3 & 4 & 12 & Alto & Mauricio González Prendas y Carlos Sainz Arce & Mitigar & Priorizar primero los módulos obligatorios y dejar mejoras visuales no esenciales para el final & Abierto \\
R-05 & Recursos humanos & Disponibilidad limitada del equipo & El equipo es pequeño y cada integrante tiene varios entregables asignados & Sobrecarga de trabajo, atrasos y poca revisión entre compañeros & 4 & 4 & 16 & Muy alto & Tamara Robles Camacho & Mitigar & Coordinar revisiones cortas, definir prioridades y reasignar tareas cuando sea necesario & Abierto \\
R-06 & Calidad & Errores críticos de navegación en el prototipo & Falta de pruebas manuales suficientes o cambios rápidos en pantallas & El usuario no podría entrar a módulos principales o continuar el flujo de navegación & 3 & 5 & 15 & Muy alto & Mauricio González Prendas y Carlos Sainz Arce & Mitigar & Ejecutar lista de verificación de navegación en inicio de sesión, panel de control y portales principales antes de la entrega & Abierto \\
R-07 & Usabilidad & Baja usabilidad del prototipo & El prototipo puede tener pantallas completas pero poco claras para el usuario & Dificultad para entender el flujo de uso por parte de estudiantes, docentes o administrativos & 3 & 4 & 12 & Alto & Mauricio González Prendas y Carlos Sainz Arce & Mitigar & Revisar claridad de menús, botones, textos y distribución visual antes del cierre & Abierto \\
R-09 & Herramientas & Problemas con el repositorio o pérdida de versiones & Uso incorrecto del repositorio, archivos mal subidos o falta de entregas de código claras & Pérdida de trazabilidad del prototipo o documentos finales & 4 & 3 & 12 & Alto & Mauricio González Prendas y Carlos Sainz Arce & Mitigar & Usar nombres claros, descripciones precisas de cambios y respaldos locales de documentos y código & Abierto \\
R-10 & Cambios & Cambios tardíos cerca de la entrega & Interesados o integrantes pueden solicitar ajustes cuando el proyecto ya está avanzado & Afectación del cronograma, calidad y revisión final & 4 & 3 & 12 & Alto & Tamara Robles Camacho, Isaac Jiménez Blanco y Jorge Abarca Quiros & Mitigar & Definir fecha límite para cambios importantes y evaluar toda solicitud mediante control de cambios & Abierto \\
R-11 & Calidad documental & Falta de criterios claros de aceptación & Si los criterios no se revisan, el equipo puede considerar listo un módulo incompleto & Entregables aceptados internamente sin cumplir navegación, contenido mínimo o coherencia & 3 & 4 & 12 & Alto & Isaac Jiménez Blanco y Jorge Abarca Quiros & Mitigar & Usar los criterios del Plan de Calidad y revisar cada módulo contra el alcance aprobado & Abierto \\
R-12 & Comunicación & Mala comunicación entre integrantes & Uso excesivo de mensajes informales sin registrar acuerdos importantes & Duplicidad de tareas, pérdida de decisiones o documentos incongruentes & 3 & 4 & 12 & Alto & Tamara Robles Camacho & Mitigar & Usar WhatsApp para avisos rápidos, pero documentar decisiones importantes en GitHub, correo o minuta & Abierto \\
R-14 & Interesados & Baja validación de interesados externos & Por tiempo limitado puede no consultarse a todos los interesados externos & El prototipo podría representar de forma incompleta algunos procesos docentes, financieros o administrativos & 3 & 4 & 12 & Alto & Tamara Robles Camacho & Aceptar y mitigar & Priorizar validaciones internas y consultar a los interesados externos solo cuando sea estrictamente necesario & Abierto \\
R-15 & Alcance técnico & Expectativa de funcionalidades no incluidas & Puede asumirse que el sistema tendrá lógica de servidor, datos reales, autenticación o pagos reales & Confusión sobre lo que realmente será entregado & 3 & 3 & 9 & Medio & Isaac Jiménez Blanco y Jorge Abarca Quiros & Evitar & Repetir en documentos y revisiones que el entregable técnico es un prototipo de interfaz navegable con datos simulados & Abierto \\
R-16 & Evidencia & Falta de evidencia final organizada & Los archivos pueden quedar dispersos entre WhatsApp, equipo local y repositorio & Dificultad para demostrar el cumplimiento de entregables y cronograma & 3 & 3 & 9 & Medio & Isaac Jiménez Blanco, Jorge Abarca Quiros, Mauricio González Prendas y Carlos Sainz Arce & Mitigar & Organizar documentos, PDFs, capturas, .mpp y código en GitHub con nombres claros & Abierto \\
\bottomrule
\end{longtable}
}
\end{widepage}

\section{Resumen de riesgos por nivel}

A partir de la evaluación de criticidad realizada, el volumen de riesgos identificados se distribuye en los términos que se exponen a continuación.

\begin{table}[H]
\centering
\caption{Resumen de riesgos por nivel}
\label{tab-resumen-riesgos}
{\small
\begin{tabularx}{\textwidth}{p{3.5cm} p{3.2cm} Y}
\toprule
\textbf{Clasificación} & \textbf{Cantidad de riesgos} & \textbf{Códigos} \\
\midrule
Crítico & 2 & R-01, R-02 \\
Muy alto & 2 & R-05, R-06 \\
Alto & 7 & R-04, R-07, R-09, R-10, R-11, R-12, R-14 \\
Medio & 2 & R-15, R-16 \\
Bajo & 0 & Ninguno \\
Muy bajo & 0 & Ninguno \\
\bottomrule
\end{tabularx}
}
\end{table}

Los riesgos clasificados en la categoría de críticos requieren una atención prioritaria y constante debido a que poseen el potencial de afectar de manera directa la entrega final de los productos. En consecuencia, estas amenazas corresponden principalmente a la restricción severa de los plazos de ejecución y al fenómeno del alcance creciente, puesto que ambos factores podrían alterar el cronograma maestro, incidir negativamente en la calidad del desarrollo y comprometer el alcance aprobado.

Asimismo, los riesgos identificados como muy altos deben ser objeto de un monitoreo estructurado ya que se asocian de manera directa con la consistencia global de la documentación, la disponibilidad horaria del equipo de trabajo y la eventual ocurrencia de fallos de navegación en el prototipo, de modo que, a pesar de no ubicarse en el extremo superior de la escala, podrían demandar retrabajo y retrasar entregables clave si no son prevenidos oportunamente.

Por otra parte, las amenazas de nivel alto se vinculan en gran medida con el diseño de la interfaz de usuario, la usabilidad de las pantallas principales, el funcionamiento de las herramientas de versionamiento, la comunicación del grupo, el procesamiento de cambios, los criterios de aceptación formales y la validación con interesados, además de los riesgos catalogados como medios que, si bien no revisten una urgencia inmediata, deben mantenerse bajo observación periódica con el fin de evitar incrementos indeseados en su nivel de severidad.

\section{Seguimiento y actualización de riesgos}

El presente registro de riesgos será sometido a actualizaciones periódicas a lo largo del ciclo de vida del proyecto en la medida en que varíen las condiciones de probabilidad, impacto, responsabilidades, acciones de respuesta o estado general de las amenazas identificadas, asimismo, se prevé la incorporación inmediata de cualquier nuevo evento detectado o la reclasificación de aquellos que se materialicen como problemas de ejecución.

La labor de seguimiento y control de estas amenazas se estructurará a través de un conjunto de actividades constantes.

\begin{itemize}[leftmargin=*, label={\textbullet}]
    \item Revisar los riesgos críticos durante las reuniones de seguimiento del equipo.
    \item Comunicar por mensajería instantánea cualquier riesgo urgente que pueda afectar la entrega.
    \item Documentar formalmente los riesgos que afecten alcance, tiempo, costo, calidad o cambios.
    \item Actualizar el registro cuando un riesgo aumente o disminuya su nivel.
    \item Relacionar los riesgos con el proceso de control de cambios cuando se requiera modificar alcance, cronograma o entregables.
    \item Revisar los riesgos asociados a la interfaz antes de la entrega final.
    \item Verificar que las acciones de respuesta se hayan aplicado cuando corresponda.
    \item Mantener evidencia organizada en el repositorio para reducir riesgos de pérdida de información.
\end{itemize}

\subsection{Estados utilizados para los riesgos}

\begin{table}[H]
\centering
\caption{Estados utilizados para los riesgos}
\label{tab-estados-riesgos}
{\small
\begin{tabularx}{\textwidth}{p{4cm} Y}
\toprule
\textbf{Estado} & \textbf{Descripción} \\
\midrule
Abierto & El riesgo fue identificado y debe mantenerse en seguimiento \\
En seguimiento & El riesgo tiene acciones activas de control \\
Materializado & El riesgo ocurrió y debe tratarse como problema \\
Cerrado & El riesgo ya no representa amenaza significativa para el proyecto \\
\bottomrule
\end{tabularx}
}
\end{table}

\subsection{Responsables de seguimiento}

\begin{table}[H]
\centering
\caption{Responsables de seguimiento}
\label{tab-responsables-seguimiento}
{\small
\begin{tabularx}{\textwidth}{p{4.5cm} Y}
\toprule
\textbf{Responsable} & \textbf{Riesgos principales asociados} \\
\midrule
Tamara Robles Camacho & Cronograma, disponibilidad del equipo, comunicación, cambios y relación con los interesados \\
Isaac Jiménez Blanco & Documentación, alcance, calidad documental, presupuesto, riesgos, cambios y evidencia \\
Jorge Abarca Quiros & Elaboración de entregables documentales, apoyo en gestión de riesgos, costos y comunicaciones del proyecto \\
Mauricio González Prendas & Desarrollo de interfaz, control de calidad, navegación, usabilidad, repositorios y correcciones técnicas \\
Carlos Sainz Arce & Desarrollador de interfaz 2, control de calidad, pruebas de calidad y navegación del portal financiero y ejecutivo \\
\bottomrule
\end{tabularx}
}
\end{table}

\section{Conclusión}

El registro de riesgos de la Plataforma Universitaria Integrada ha facilitado la oportuna identificación y priorización de las diversas amenazas que podrían obstaculizar el cumplimiento de los objetivos de la organización, de manera que este instrumento asiste activamente al equipo en la anticipación de desviaciones asociadas a los plazos, el alcance técnico, la calidad del producto, los procesos de comunicación, el uso de plataformas colaborativas y las actividades finales de cierre.

Las amenazas de mayor relevancia se concentran principalmente en la restricción severa de tiempo y en la posibilidad de alcance creciente, además de la consistencia de los documentos, la disponibilidad de los colaboradores, el correcto funcionamiento de la navegación y la estimación del esfuerzo de desarrollo de la interfaz de usuario. Asimismo, como resultado de la evaluación cuantitativa, la limitación temporal y el crecimiento descontrolado del alcance se definieron como riesgos críticos, mientras que la coherencia documental, la disponibilidad de personal y los posibles fallos del prototipo se catalogaron como de muy alta severidad, por lo tanto, estos eventos recibirán un seguimiento sumamente estrecho en la fase de cierre.

Este registro favorece la asignación de responsabilidades claras para el tratamiento de cada evento identificado, de modo que la dirección del proyecto coordinará el seguimiento de los riesgos de cronograma, comunicación y relaciones con interesados, por su parte, los encargados del área de análisis y documentación velarán por la gestión del alcance, los presupuestos, la trazabilidad documental y los cambios formales, asimismo, los desarrolladores y evaluadores técnicos controlarán de forma permanente la navegación, el control de versiones, la usabilidad y la corrección de errores en la interfaz.

Mediante la implementación continua de este registro, el equipo de trabajo dispondrá de un marco estructurado para monitorear las amenazas, ejecutar planes de mitigación preventivos y actuar de forma coordinada en caso de que alguna contingencia llegue a materializarse. En consecuencia, esta gestión dinámica contribuirá a resguardar la integridad del alcance técnico, mantener la planificación temporal dentro de los rangos tolerables, asegurar la calidad final de los productos y garantizar la alineación del proyecto con la totalidad de los planes y documentos estratégicos establecidos previamente.

\end{document}
